\section{Positive forms on fusion products}
\label{sec:fusion}

We now use the normal representations constructed in
Theorem~\ref{thm:integration} to transfer a vacuum Gram matrix to an
arbitrary unitary module.  The algebraic fusion product need not yet be
unitary.  We first construct a nondegenerate invariant sesquilinear form
on that product, and then prove its positivity separately on every
conformal-energy space.

For unitary modules $M,N$, extend the diagonal transport form
\eqref{eq:transport-form} to positive real insertion points by setting
\begin{equation}
\begin{split}
 \mathfrak t^{M,N}_{z,w}(a\otimes b,a'\otimes b')
 =\big\langle
 Y^N\big(\mathcal Y_M^-(\theta_w a',w^{-1}-z)a,z\big)b,b'
 \big\rangle_N,
\end{split}
\label{eq:two-point-transport}
\end{equation}
where $a,a'\in M$, $b,b'\in N$, and
\begin{equation}
 0<z,w<1,\qquad w^{-1}>z>w^{-1}-z>0.
 \label{eq:positive-transport-region}
\end{equation}
As usual, insertion of a vector in the grade completion means the
absolutely convergent sum over its homogeneous components.  Thus
$\mathfrak t^{M,N}_{r,r}=\mathfrak t^{M,N}_r$ in the notation of
Section~\ref{sec:annuli}.

\begin{theorem}[Positive fusion forms]
\label{thm:positive-fusion}
Let $V$ be a simple unitary strongly rational VOA.  For any unitary
grading-restricted finite-length $V$-modules $M,N$, let
$X=M\boxtimes N$ and let $\mathcal Y_X$ be its universal intertwining
operator.  There is a positive-definite invariant Hermitian form $h_X$
on $X$ such that
\begin{equation}
 \mathfrak t^{M,N}_{z,w}(a\otimes b,a'\otimes b')
 =\sum_\lambda h_X\big(
 \Pi_\lambda\mathcal Y_X(a,z)b,
 \Pi_\lambda\mathcal Y_X(a',w)b'\big)
 \label{eq:positive-fusion-identity}
\end{equation}
for all the vectors and coordinates in
\eqref{eq:two-point-transport}--\eqref{eq:positive-transport-region}.
The sum converges absolutely.  This is the canonical positive
normalization of the fusion form.
\end{theorem}

We prove the theorem in the next three subsections and check its
normalization in the fourth.

\subsection{A positive matrix in the local algebra}

We begin with positivity of the diagonal form.  Fix an interval $I$ near
$1\in S^1$, a localized dented annulus $\underline A$ as in
Section~\ref{sec:annuli}, and its double
$\widehat A=\underline A^\dagger\underline A$.  Choose a real $r<1$
sufficiently close to $1$ in its inner collar.  This choice works for
every finite list of charge vectors; constants in estimates may depend
on that list.  Write $\delta=r^{-1}-r$ and $q=2\delta$.

For $a_1,\ldots,a_m\in M$, set
\[
 u_{ij}=\mathcal Y_M^-(\theta_r a_i,\delta)a_j
 \quad\text{in the grade completion of }V.
\]
By linearity and \eqref{eq:annulus-hilbert-state},
$q^{-L_0}u_{ij}\in\mathcal H_V$.  The fixed-disc estimate
\eqref{eq:fixed-disc-bound} therefore defines bounded operators
\begin{equation}
 K_{ij}=\widehat A_V[u_{ij}(r)]\in A_V(I)
 \label{eq:local-transport-matrix}
\end{equation}
as operator-norm limits of the insertions of the finite level
truncations of $u_{ij}$.  In particular these limits remain in the
same local von Neumann algebra.

The vacuum transport identity for creation, namely
Proposition~3.10 of~\cite{TenerFusion}, and the insertion identity
\eqref{eq:annulus-insertion} give
\begin{equation}
 K_{ij}=T_i(r)^*T_j(r),\qquad
 T_j(r)=\mathcal Y_M^+(a_j,r)\underline A_V.
 \label{eq:transport-gram}
\end{equation}
Here the creation operators with the displayed cutoff are bounded by
Lemma~\ref{lem:charged-tools}.  The equality is first an equality of
matrix elements on $\underline A_V^{-1}V$, which is dense by
Lemma~\ref{lem:annulus-range}; boundedness extends it to the Hilbert
space.  Thus $(K_{ij})_{i,j=1}^m$ is a positive element of
$M_m(A_V(I))$.

By Theorem~\ref{thm:integration}, $N$ has a normal local representation
$\pi_I^N$.  Applying its matrix amplification preserves this
positivity.  For arbitrary $b_1,\ldots,b_m\in N$, choose exact
preimages $\xi_i$ under $\underline A_N$ as follows.  If the
complementary annulus of Lemma~\ref{lem:annulus-range} satisfies
$\underline A_N\underline B_N=s^{L_0}$, put
\[
 \xi_i=\underline B_Ns^{-L_0}b_i.
\]
The inverse power acts on a finite-energy vector, so this expression
is defined and $\underline A_N\xi_i=b_i$.  The norm-convergent
insertion identity, followed by absolute convergence of the
transport series, now gives
\begin{equation}
 \sum_{i,j=1}^m
 \mathfrak t^{M,N}_{r,r}(a_j\otimes b_j,a_i\otimes b_i)
 =\sum_{i,j=1}^m
 \langle\pi_I^N(K_{ij})\xi_j,\xi_i\rangle_N
 \geq0.
 \label{eq:transport-positive}
\end{equation}
Only the creation case was used to obtain the vacuum matrix.  General
fusion positivity has not entered this argument.

\subsection{The invariant form on the algebraic fusion object}

Assume first that $M,N$ are simple.  Associativity in the
positive-real region produces an intertwining operator $B$ of type
$\binom{N}{\bar M\ X}$ such that
\begin{equation}
 Y^N\big(\mathcal Y_M^-(c,w^{-1}-z)a,z\big)b
 =B(c,w^{-1})\mathcal Y_X(a,z)b
 \label{eq:fused-evaluation}
\end{equation}
in convergent matrix elements, for $c\in\bar M$, $a\in M$ and
$b\in N$.  Its morphism is the associator
followed by the annihilation morphism
$\bar M\boxtimes M\longrightarrow V$ and the left unit action on
$N$.  The annihilation is nonzero: its adjoint is creation, which
sends $a\otimes\Omega$ to $e^{zL_{-1}}a$.  Since $M$ is simple,
unitarity identifies $\bar M$ with $M^\vee$, and the annihilation
morphism is a nonzero multiple of rigid evaluation.  We also write
$B$ for its corresponding morphism.  Consequently
rigidity gives an isomorphism
\begin{equation}
 \operatorname{Hom}_V(E,X)\longrightarrow
 \operatorname{Hom}_V(\bar M\boxtimes E,N),
 \qquad f\longmapsto B(1_{\bar M}\boxtimes f)
 \label{eq:rigid-transport-adjunction}
\end{equation}
for every grading-restricted finite-length module $E$.

To turn \eqref{eq:fused-evaluation} into a form, recall the Hermitian
duality from Section~\ref{sec:preliminaries}, which does not require
positive forms.  For a module $E$, its conjugate
$\bar E$ has mode action
\[
 v_{(n)}^{\bar E}\bar x
   =\overline{(\Theta v)_{(n)}^E x}.
\]
Its Hermitian dual is $E^{\mathrm h}=(\bar E)^\vee$, with pairing
\begin{equation}
 [x,\lambda]_E=\overline{\lambda(\bar x)}.
 \label{eq:hermitian-dual-pairing}
\end{equation}
The pairing is linear in $x$ and conjugate-linear in $\lambda$.
Restricted contragredience gives its invariant adjoint relation.
The given inner product of $N$ is represented by a linear module
isomorphism $H_N:N\longrightarrow N^{\mathrm h}$.

Conjugate $B$ and apply the FHL contragredient transformation
\cite[Section~5.5]{FHL}, in the explicit algebraic form
\cite[Equation~(3.87), Proposition~3.46]{HLZII}, taking the integer
branch index in that formula to be $-1$.  Thus we use
\[
 e^{wL_1}e^{-i\pi L_0}w^{-2L_0}
\]
on the conjugated charge, and precompose the second input with
$H_N$.  This gives an intertwining
operator $B^\sharp$ of type $\binom{X^{\mathrm h}}{M\ N}$ with
\begin{equation}
 [x,B^\sharp(a,w)b]_X
   =\langle B(\theta_w a,w^{-1})x,b\rangle_N
 \qquad(x\in X,\ a\in M,\ b\in N, w>0).
 \label{eq:hermitian-dual-adjoint}
\end{equation}
The minus phase before conjugation gives the plus phase in
$\theta_w$.  All module weights are real, so these transformations
also preserve the prescribed positive-real branches.
Universality of $\mathcal Y_X$ gives a module map
$H:X\longrightarrow X^{\mathrm h}$ such that
$B^\sharp=H\mathcal Y_X$.

The map $H$ is surjective.  To see the role of rigidity, suppose that
the quotient map $q:X^{\mathrm h}\longrightarrow C=\operatorname{coker}H$
has nonzero target.  Hermitian duality identifies
$(X^{\mathrm h})^{\mathrm h}$ with $X$, so dualizing $q$ gives a
nonzero injection
\[
 f=q^{\mathrm h}:E=C^{\mathrm h}\longrightarrow X,
 \qquad f^{\mathrm h}H=qH=0.
\]
Pairing \eqref{eq:hermitian-dual-adjoint} with $f(E)$ would imply
$B(1_{\bar M}\boxtimes f)=0$: the charge transformation $\theta_w$
is invertible, and the pairing with $N$ is nondegenerate.  This
contradicts \eqref{eq:rigid-transport-adjunction}.  Hermitian duality
preserves finite length, so the surjective map $H$ is an isomorphism.

We have therefore constructed the nondegenerate invariant
sesquilinear form
\begin{equation}
 h_X(x,y)=[x,Hy]_X.
 \label{eq:algebraic-fusion-form}
\end{equation}
It pairs distinct energy spaces to zero.  At this stage neither
Hermitian symmetry nor positivity has been assumed.  Combining
\eqref{eq:fused-evaluation} and
\eqref{eq:hermitian-dual-adjoint} proves
\eqref{eq:positive-fusion-identity} for this form.  Absolute
convergence is the ordinary convergence of the product of $B$ and
$\mathcal Y_X$, applied to the restricted-dual functional supplied
by a finite-energy vector of $N$.

\subsection{Extracting positivity at one energy}

The positivity in \eqref{eq:transport-positive} concerns a sum over
all output energies.  The following lemma separates those energies
without first constructing a norm on $X$.

\begin{lemma}[Positive polynomial moments]
\label{lem:positive-moments}
Let $\Lambda\subset\mathbb R$ be countable and bounded below.
Suppose complex numbers $c_\lambda$ satisfy
\[
 \sum_{\lambda\in\Lambda}e^{\epsilon\lambda}|c_\lambda|<\infty
 \quad\text{for some }\epsilon>0.
\]
If $\sum_\lambda |p(\lambda)|^2c_\lambda$ is real and nonnegative
for every complex polynomial $p$, then every $c_\lambda$ is real
and nonnegative.
\end{lemma}

\begin{proof}
Set $\mu=\sum_\lambda c_\lambda\delta_\lambda$.  Polynomials are
dense in $L^2(|\mu|)$.  To see this, let $g$ be orthogonal to them.
Cauchy--Schwarz and the exponential moment show that the finite
complex measure $g\,|\mu|$ has an exponential moment of smaller
positive order.  Its bilateral Laplace transform is holomorphic in
a strip containing the imaginary axis.  Orthogonality makes every
derivative at zero vanish, so the identity theorem makes the
transform zero on this strip.  Uniqueness of the Fourier transform
of a finite measure then gives $g\,|\mu|=0$.  For completeness, this
last uniqueness follows by convolving the measure with Gaussian
approximate identities: every convolution has zero Fourier
transform and is zero, and the convolutions converge weakly to the
measure.

For a fixed $\lambda$, approximate the indicator of its atom by
polynomials $p_n$ in $L^2(|\mu|)$.  Then
$|p_n|^2$ converges to that indicator in $L^1(|\mu|)$, because
\[
 \big\||p_n|^2-|p_m|^2\big\|_{L^1(|\mu|)}
 \leq\|p_n-p_m\|_{L^2(|\mu|)}
       \big(\|p_n\|_{L^2(|\mu|)}+\|p_m\|_{L^2(|\mu|)}\big).
\]
Thus the nonnegative integrals $\int |p_n|^2\,d\mu$ tend to
$c_\lambda$.
\end{proof}

Fix the radius $r$ used in \eqref{eq:transport-positive}.  Consider
a finite sum in the grade completion,
\[
 x=\sum_i\mathcal Y_X(a_i,r)b_i,
 \qquad x_\lambda=\Pi_\lambda x,
 \qquad c_\lambda=h_X(x_\lambda,x_\lambda).
\]
After decomposing the inputs, we may take all $a_i,b_i$
homogeneous.  The numbers $c_\lambda$ have an exponential moment
in absolute value.  Indeed, choose $z>r$ sufficiently close to $r$
that $z<r^{-1}$.  In the absolutely convergent product with the
outer insertion fixed at $r^{-1}$, replacing the inner radius $r$
by $z$ multiplies its component of weight $\lambda$ by
\[
 (z/r)^{\lambda-\operatorname{wt}(a_i)-\operatorname{wt}(b_i)}.
\]
Absolute product convergence for each of the finitely many pairs
of inputs gives
$\sum_\lambda e^{\epsilon\lambda}|c_\lambda|<\infty$ for a
common $\epsilon>0$.  This uses only the sesquilinear form
\eqref{eq:algebraic-fusion-form} and algebraic convergence.

The space of such finite sums $x$ is stable under $L_0$, since
\begin{equation}
 L_0\mathcal Y_X(a,r)b
 =\mathcal Y_X(L_0a,r)b+\mathcal Y_X(a,r)L_0b
   +r\mathcal Y_X(L_{-1}a,r)b.
 \label{eq:fixed-radius-L0}
\end{equation}
The radius $r$ stays fixed in this step.  Each polynomial in $L_0$
introduces only finitely many further charge and source vectors, and
the fixed-disc construction of \eqref{eq:transport-positive} applies
at this same $r$ to every such finite list.  Its constants may change;
its positivity conclusion does not.
Apply \eqref{eq:transport-positive} to $p(L_0)x$.  Grade
orthogonality gives
\[
 \sum_\lambda |p(\lambda)|^2h_X(x_\lambda,x_\lambda)\geq0.
\]
Lemma~\ref{lem:positive-moments} proves
$h_X(x_\lambda,x_\lambda)\geq0$ at each energy.

These $x_\lambda$ run through the whole space $X_\lambda$.  In
fact the span of all coefficients of $\mathcal Y_X(a,z)b$ is a
graded submodule by the intertwiner commutator formula.  The
quotient map by that submodule has zero postcomposition with the
universal intertwiner, and is therefore zero.  For homogeneous
inputs each projection at the fixed nonzero radius $r$ is a
nonzero scalar multiple of its coefficient.  This proves the
asserted spanning at each energy.

It follows by polarization that $h_X$ is Hermitian and positive
semidefinite on every energy space, and hence on $X$.  A
zero-norm vector for a positive semidefinite Hermitian form belongs
to its radical.  Nondegeneracy of \eqref{eq:algebraic-fusion-form}
therefore makes $h_X$ positive definite.  In particular, $X$ is now
a unitary module.  Finite orthogonal decompositions of $M,N$
extend the construction to arbitrary unitary finite-length
modules: creation, annihilation and the transport identity split
in exactly these orthogonal blocks.

\subsection{Canonical normalization and general charged fields}

With $h_X$ now positive, \eqref{eq:hermitian-dual-adjoint} says
that $B$ is the adjoint of $\mathcal Y_X$ in Tener's convention.
Thus \eqref{eq:fused-evaluation} is the exact normalized
fusion identity for that convention.

We check the normalization, since fractional charge weights make
the adjoint phase significant.  Put, for $x>0$,
\[
 U_\pm(x)=e^{xL_1}e^{\pm i\pi L_0}x^{-2L_0}.
\]
Gui's contragredient convention is
$U_+(x)$ in~\cite[Equations~(8.3)--(8.5)]{GuiPositivity}; solving
for the adjoint uses
$U_-(x)=U_+(x^{-1})^{-1}$.  On a simple charge module of lowest
weight $h$, $U_-=e^{-2\pi i h}U_+$.  With inner products linear in
the first variable, this means that Gui's adjoint, after the
unitary identification of the conjugate charge with the
contragredient charge, is $e^{2\pi i h}$ times Tener's adjoint.
Gui's annihilation operator has the same factor relative to
$\mathcal Y_M^-$: its vacuum coefficient on a top vector $a$ is
$e^{i\pi h}x^{-2h}\|a\|^2$, whereas the coefficient for
$\mathcal Y_M^-$ is $e^{-i\pi h}x^{-2h}\|a\|^2$.

The two factors cancel in the fusion identity
\cite[Equation~(8.8)]{GuiPositivity}.  Thus the canonical operator
called $A$ in that identity is $1$ for the form just constructed.
The identity leaves no scalar freedom: the coefficients of the
universal intertwiner span $X$, and their pairings determine its form.
This proves the canonical-normalization assertion in
Theorem~\ref{thm:positive-fusion}.  The comparison applies on each
orthogonal simple charge summand, so also covers finite sums.
Common-angle versions follow by simultaneous $L_0$ rotations,
using conjugate logarithms at reflected coordinates.  In
particular \eqref{eq:positive-fusion-identity} at a rotated
diagonal is the squared norm of the universal intertwiner there.

We have proved closure of the unitary modules under the full
algebraic fusion product, without yet knowing that every simple
module is unitary.  The following consequence supplies the
charged fields needed for the comparison with conformal nets.

\begin{corollary}
\label{cor:charged-general}
Let $V$ be a simple unitary strongly rational VOA, and let $M,N,K$
be unitary grading-restricted finite-length $V$-modules.  Every
intertwining operator $\Phi$ of type $\binom{K}{M\ N}$ is
polynomially energy-bounded for every homogeneous charge vector
and has the strong intertwining property with all observable
fields.  More precisely, for homogeneous $a\in M$ of weight $d$,
write
\[
 \Phi(a,z)=\sum_\mu\Phi(a)[\mu]z^{-\mu-d}.
\]
There are $C>0$ and nonnegative integers $s,k$ such that
\[
 \|\Phi(a)[\mu]b\|
 \leq C(1+|\mu|)^s\|(1+L_0)^k b\|
 \qquad(b\in N).
\]
If a smooth test $f$ is supported in a lifted interval $I$, the
closed off-diagonal smear $\overline{\Phi(a,f)}$ on
$\mathcal H_N\oplus\mathcal H_K$ is affiliated with the
commutant of the algebra generated by
\[
 \operatorname{diag}\big(\overline{Y^N(v,g)},
                         \overline{Y^K(v,g)}\big),
 \qquad v\in V,\quad \operatorname{supp}g\subset J,
\]
for every interval $J$ disjoint from $I$.
\end{corollary}

\begin{proof}
By universality, $\Phi=L\mathcal Y_X$ for a module map
$L:X=M\boxtimes N\longrightarrow K$.  Both modules are now
unitary.  Such a map is bounded: decompose into the finitely many
unitary simple summands; Schur's lemma and uniqueness of an
invariant form on each simple identify it with a finite family of
linear maps between finite-dimensional multiplicity spaces.
Theorem~\ref{thm:positive-fusion}, including its rotated-diagonal
identity, gives
\[
 \|\Phi(a,z)b\|^2
 \leq\|L\|^2\,
       \mathfrak t_z^{M,N}(a\otimes b,a\otimes b).
\]
The source $N$ is strongly integrable by
Theorem~\ref{thm:integration}.  This is precisely the transport
domination hypothesis of Lemma~\ref{lem:charged-tools}, which
gives the stated mode bounds and bounded annulus-cutoff smears.
Proposition~\ref{prop:strong-intertwining} then gives the asserted
affiliation and strong intertwining property for homogeneous
observables.  For a finite sum of homogeneous observables, sum the
weak commutation identities on the common smooth core before
closing.  The same bounded regularization and its adjoint preserve
that core, so the graph argument of that proposition also applies
to the closure of the sum.  This proves the assertion for every
$v\in V$.
\end{proof}
